\section{重回帰分析と変数間の影響}
\label{b08_MultipleRegression}

\subsection{授業内容}
\subsubsection{概要}
前回は単回帰分析について学んだ。今回は複数の説明変数を用いる重回帰分析へと発展させる。回帰分析によって，データにモデルを当てはめて係数の推定値を得るということを行ったが，そのこととモデルが適していたかどうかは別である。この点を理解するためには，モデルの適合度について考える必要がある。モデルが適合しているということを踏まえた上で，説明する変数が増える重回帰分析へと議論は展開する。ここでは説明変数の読み取り方への注意を促すべく，偏相関係数の考え方を解説し，偏回帰係数，標準化偏回帰係数などの用語について理解する。

\subsubsection{コマ主題細目}
\begin{description}
	\item[重回帰分析の基本] 説明変数が複数に増えた場合の回帰分析は，重回帰分析(Multiple Regression Analysis)と呼ぶ。数式的な表現としては項が増えるだけであるが，データのプロットを考えると三次元以上の回帰平面に拡張されていること，それでも面は湾曲していない，すなわち線形関係であることを理解する。単回帰分析と同様に，最小二乗法によって回帰係数を推定するが，行列計算を用いることで効率的に解が求められることにも触れる。
	      \begin{flushright}
		      $\to$ \textcite{kosugi2019}のP.168,図10.7の回帰平面図を参照
	      \end{flushright}


	\item[部分相関と偏相関と偏回帰係数] 重回帰分析における係数の解釈のために，部分相関と偏相関の概念を理解する。回帰分析によって，被説明変数の分散が説明される部分と残差の部分に分離させられること，残差と説明変数は相関しないことを確認し，残差は説明変数の影響を除外したものと考えることができる。これは条件を必ずしも統制できない事象において，統計モデル的に統制していることでもある。偏相関係数は他の変数の影響を取り除いた2変数間の純粋な関連の強さを表す指標であり，変数間の直接的な関係を評価する上で重要な概念である。
	      \begin{flushright}
		      $\to$ \textcite{Kawabata201408}のP.59-65，\textcite{Yamada2004}のP.60-65。
	      \end{flushright}

	\item[偏回帰係数と解釈] 部分相関の説明において，残差が分離されることを見てきたが，残差同士の相関関係はとくに偏相関ということができる。これは共通する変数からの影響を除いた変数間関係であり，これを見るだけでも変数間関係を見ることにつながる。とくに残差変数を使っての回帰分析を行うとき，得られる係数が偏回帰係数である。重回帰分析における回帰係数はこのように条件付きで考えなければならない。偏回帰係数は「他の説明変数の値を一定に保った場合に，その説明変数が1単位変化したときの被説明変数の変化量」として解釈できることを理解する。また変数を増やすことで重相関係数が増加することも理解する。また，複数の説明変数間に強い相関がある場合，多重共線性(multicollinearity)の問題が生じる可能性があることを理解する。多重共線性があると，回帰係数の推定値が不安定になり，解釈が困難になるという問題が発生する。VIF(Variance Inflation Factor)などの診断指標や，多重共線性が検出された場合の対処法についても簡単に触れる。
	      \begin{flushright}
		      $\to$\textcite{kosugi2018}のP.60--63.多重共線性については\textcite{Shimizu201705}のP.86--89を参照。
	      \end{flushright}

	\item[標準化された係数] 回帰係数については，素点による回帰係数が理解しやすいが，単位に左右されるものでもあるので相対比較をする場合はすべての変数を標準化した標準化解を用いた方がわかりやすい。標準化解を用いても適合度は変化しないことなど留意点とともに有用性を理解する。
	      \begin{flushright}
		      $\to$\textcite{kosugi2018}のP.66--67，
	      \end{flushright}

\end{description}
\subsubsection{キーワード}
\begin{itemize}
	\item 重回帰分析と回帰平面
	\item 交絡要因と統計的制御
	\item 部分相関と偏相関
	\item 偏回帰係数と標準化偏回帰係数
	\item 多重共線性
\end{itemize}
\subsection{授業情報}
\paragraph{コマの展開方法}
講義と演習
\paragraph{標準シラバスにおける位置づけ}
科目番号5 心理学統計法；1.心理学で用いられる統計手法；C 記述統計：回帰
\subsubsection{予習・復習課題}
\paragraph{予習}
単回帰分析について，数式的表現や最小二乗法の意味について概念的な理解で十分なので再確認しておくこと。重回帰分析では説明変数が複数になるため，解析的算出には行列計算の知識が必要になる。行列計算については，二年次の統計に関する講義で詳しく扱う予定であるが，興味があるものは「線形代数」をキーワードに予習すると良い。初学者には\textcite{Yuki201810}または\textcite{Okada200804}が適している。
\paragraph{復習}
\textcite{Toyoda201706}のP.89--98は本時の内容がわかりやすくまとめられているので，通読しておくと良い復習になる。また，授業で学んだ内容を応用して以下の課題に取り組むこと：

\begin{enumerate}
\item 提供されたデータセットを用いて，Rで重回帰分析を実施する
\item 各説明変数の偏回帰係数と標準化偏回帰係数を求め，その意味を解釈する
\item モデルの決定係数と調整済み決定係数を算出し，モデルの適合度を評価する
\item 多重共線性の診断を行い，問題がある場合は対処法を考察する
\item 交絡要因の影響を統制した結果，単回帰分析と比較してどのような変化があったかを考察する
\end{enumerate}

重回帰分析は心理学研究で広く用いられる手法なので，\textcite{kosugi2016a}の第二章などを読んで，どのような使い方をされているのか，また前提となっている条件や利用上の問題点などについても理解を深めておくこと。